\begin{textsample}{POS Dim 1 – vem\_ed\_23 – Score 3.00 – vem\_ed\_23/t118.txt}  \label{ex:f1_pos_146}
Conferência Faubai 2024: Big Data e IA nos PCIs A equipe dos Projetos Colaborativos Internacionais (PCIs / Cesu) teve três trabalhos aprovados, além de \textbf{ser} convidada a participar da moderação de sessões e de mesa-redonda na Faubai 2024, maior e mais longeva conferência sobre internacionalização da educação \textbf{superior} da América Latina, realizada de 21 a 24 de abril no Centro de Convenções Rebouças, em São Paulo.
Temas emergentes, como o uso de técnicas de Big Data e Inteligência Artificial no planejamento, design e avaliação de projetos COIL (Collaborative Online International Learning), ou Intercâmbios Virtuais, compuseram a pauta de dois trabalhos apresentados por Osvaldo Succi Junior, coordenador da equipe dos PCIs / Cesu.
O primeiro, elaborado em coautoria com Patrícia Sales Patrício, responsável pela comunicação dos PCIs / Cesu, e Carlos Eduardo Dantas de Menezes, professor da Fatec Ipiranga, \textbf{foi} apresentado no dia 22 de abril às 10h.
Intitulada “Improving COIL coordinators’ Decision-Making Process through Data-Informed Strategies”, a comunicação oral lotou a sala.
Em tradução livre, o título \textbf{é} “Aprimoramento do \textbf{processo} de tomada de decisão dos coordenadores de COIL [ Intercâmbios Virtuais ] por meio de \textbf{estratégias} informadas por dados".
Trata-se de uma pesquisa exploratória conduzida por Menezes com seus alunos do curso de Big Data para Negócios a \textbf{partir} das pesquisas de percepção aplicadas pela equipe dos PCIs / Cesu aos alunos e professores envolvidos nos PCIs desde 2020.
Técnicas como árvore da decisão, bag of words (literalmente, “saco de palavras”) e diagrama de Venn \textbf{foram} \textbf{utilizadas} para ensaiar modelos preditivos.
A técnica da árvore da decisão, por exemplo, demonstrou \textbf{que} clareza dos objetivos de aprendizagem no design de PCIs \textbf{é} um dos principais \textbf{fatores} de sucesso.
E o \textbf{que} \textbf{seria} sucesso de um PCI?
Manutenção do projeto por mais de duas edições.
Na sessão imediatamente a \textbf{seguir} (10h20), Succi Junior falou sobre “Using Artificial Intelligence as an Assistive Tool for COIL VE Coordinators in Project Design” (traduzindo, “Uso da inteligência \textbf{artificial} como ferramenta de auxílio para coordenadores de Intercâmbios Virtuais no design de projetos”).

% matched lemmas: artificial, estratégia, fator, partir, processo, que, seguir, ser, superior, utilizar
\end{textsample}
